\documentclass[a4paper,11pt]{article}

\usepackage[utf8]{inputenc}
\usepackage[T1]{fontenc}
\usepackage{amsmath, amssymb, amsthm}
\usepackage{geometry}
\geometry{margin=2.5cm}

\title{Definizione dell'integrale di una forma su una superficie\\
tramite approssimazione diadica}
\author{}
\date{}

\theoremstyle{definition}
\newtheorem{definition}{Definizione}
\newtheorem{theorem}{Teorema}
\newtheorem{proposition}{Proposizione}

\begin{document}

\maketitle

\section{Idea generale}

Costruiamo l'integrale di una $k$-forma differenziale su una superficie parametrizzata
utilizzando approssimazioni su reticoli diadici e correnti poliedrali.

L'idea è:
\begin{itemize}
\item discretizzare lo spazio tramite una griglia diadica;
\item definire l'integrazione su oggetti discreti;
\item approssimare superfici lisce con tali oggetti;
\item passare al limite.
\end{itemize}

\section{Correnti diadiche}

Sia $\mathcal{G}_N$ la griglia diadica di passo $2^{-N}$ in $\mathbb{R}^n$.

\begin{definition}
Una \emph{$k$-corrente diadica} è una combinazione lineare finita di celle
$k$-dimensionali della griglia $\mathcal{G}_N$, orientate.
\end{definition}

Se $P$ è una corrente diadica e $\omega$ è una $k$-forma liscia, si definisce:
\[
P(\omega) := \sum_i \int_{Q_i} \omega,
\]
dove $Q_i$ sono le celle.

\section{Teorema di Stokes (caso diadico)}

\begin{proposition}
Se $P$ è una corrente diadica e $\omega$ è una $(k-1)$-forma liscia, allora:
\[
P(d\omega) = \partial P(\omega).
\]
\end{proposition}

\noindent
La dimostrazione è immediata: vale su ogni cella e le facce interne si cancellano.

\section{Approssimazione di correnti}

Sia $T$ una $k$-corrente (regolare). Esiste una successione di correnti diadiche $P_N$
tali che:
\[
T = P_N + \partial S_N + R_N,
\]
con stime:
\[
\mathbf{M}(S_N) \le C 2^{-N}, \qquad \mathbf{M}(R_N) \to 0.
\]

In particolare, per ogni forma liscia $\omega$:
\[
P_N(\omega) \to T(\omega).
\]

\section{Correnti parametrizzate}

Sia $\varphi : Q \subset \mathbb{R}^k \to \mathbb{R}^n$ una funzione continua.

\begin{definition}
La corrente parametrizzata associata a $\varphi$ è definita da:
\[
T_\varphi(\omega) := \lim_{N\to\infty} P_N(\omega),
\]
dove $P_N$ è una approssimazione diadica della superficie $\varphi(Q)$.
\end{definition}

La convergenza segue dalle stime:
\[
|(P_N - P_M)(\omega)| \le C 2^{-N} \|d\omega\|_\infty + o(1),
\]
quindi $(P_N(\omega))$ è di Cauchy.

\section{Formula classica (caso regolare)}

\begin{theorem}
Sia $\varphi : Q \subset \mathbb{R}^k \to \mathbb{R}^n$ di classe $C^1$.
Allora, per ogni $k$-forma liscia $\omega$:
\[
T_\varphi(\omega) = \int_Q \varphi^* \omega.
\]
\end{theorem}

\subsection*{Idea della dimostrazione}

\begin{itemize}
\item Si suddivide $Q$ in cubi diadici $Q_i$ di lato $2^{-N}$.
\item Su ogni cubo:
\[
\varphi(x) = \varphi(x_i) + D\varphi(x_i)(x-x_i) + O(2^{-2N}).
\]
\item Quindi $\varphi$ è approssimata da una mappa affine.
\item L'integrale su ciascun cubo è approssimato dall'integrale sul piano tangente.
\item L'errore totale è $O(2^{-N}) \to 0$.
\end{itemize}

Passando al limite si ottiene:
\[
\lim_{N\to\infty} P_N(\omega) = \int_Q \varphi^*\omega.
\]

\section{Forma esplicita}

Se
\[
\omega = f\, dx^{i_1} \wedge \cdots \wedge dx^{i_k},
\]
allora:
\[
\varphi^*\omega =
f(\varphi(u)) \cdot
\det\left( \frac{\partial(\varphi^{i_1}, \dots, \varphi^{i_k})}{\partial(u_1, \dots, u_k)} \right)
\, du_1 \cdots du_k.
\]

\section{Conclusione}

Abbiamo costruito l'integrale di una forma su una superficie:
\begin{itemize}
\item a partire da oggetti discreti (correnti diadiche);
\item passando al limite;
\item recuperando la formula classica nel caso regolare.
\end{itemize}

\section{Lemma di deformazione (versione diadica)}

\begin{theorem}[Lemma di deformazione]
Sia $T$ una $k$-corrente a massa finita in $\mathbb{R}^n$.
Per ogni $N \in \mathbb{N}$, sia $\mathcal{G}_N$ la griglia diadica di passo $2^{-N}$.

Allora esistono:
\begin{itemize}
\item una $k$-corrente diadica $P_N$ supportata sullo scheletro $k$-dimensionale di $\mathcal{G}_N$,
\item una $(k+1)$-corrente $S_N$,
\item una $k$-corrente $R_N$,
\end{itemize}
tali che:
\[
T = P_N + \partial S_N + R_N,
\]
e valgono le stime:
\[
\mathbf{M}(P_N) \le C\,\mathbf{M}(T),
\]
\[
\mathbf{M}(S_N) \le C\,2^{-N}\,\mathbf{M}(T),
\]
\[
\mathbf{M}(R_N) \le C\,2^{-N}\,\mathbf{M}(\partial T),
\]
dove $C$ dipende solo da $n$ e $k$.

Inoltre, se $\partial T = 0$, allora si può scegliere $R_N = 0$, cioè:
\[
T = P_N + \partial S_N.
\]
\end{theorem}

\medskip

\noindent
\textbf{Conseguenza.}
Per ogni $k$-forma liscia $\omega$ con coefficienti limitati,
si ha:
\[
\lim_{N\to\infty} P_N(\omega) = T(\omega).
\]

Infatti, per $N < M$:
\[
(P_M - P_N)(\omega)
= (S_N - S_M)(d\omega) + (R_N - R_M)(\omega),
\]
da cui segue la stima:
\[
|(P_M - P_N)(\omega)|
\le C\,2^{-N}\,\|d\omega\|_\infty
+ o(1),
\]
che implica che la successione $(P_N(\omega))$ è di Cauchy.

\section{Dimostrazione del lemma di deformazione}

\begin{proof}
L'idea della dimostrazione è deformare progressivamente la corrente $T$
verso lo scheletro $k$-dimensionale della griglia diadica $\mathcal{G}_N$,
riducendo passo dopo passo la dimensione del supporto.

\medskip

\noindent
\textbf{Passo 1: riduzione allo scheletro $(n-1)$-dimensionale.}

Sia $\mathcal{G}_N$ la griglia diadica di passo $2^{-N}$.
Consideriamo un cubo $Q$ di dimensione $n$ della griglia.

Fissato un punto $x_Q$ nell'interno di $Q$, definiamo una mappa di proiezione radiale
\[
\pi_Q : Q \setminus \{x_Q\} \to \partial Q
\]
che manda ogni punto sul bordo lungo il segmento che unisce $x_Q$ al punto stesso.

Questa mappa è lipschitziana su ogni compatto di $Q \setminus \{x_Q\}$.
Applicando $\pi_Q$ a ciascun cubo e incollando le mappe,
otteniamo una deformazione globale
\[
\pi^{(n)} : \mathbb{R}^n \to \mathbb{R}^n
\]
che manda lo spazio nello scheletro $(n-1)$-dimensionale della griglia.

Definiamo la corrente deformata:
\[
T^{(n-1)} := \pi^{(n)}_{\#} T.
\]

\medskip

\noindent
\textbf{Passo 2: formula di omotopia.}

La mappa $\pi^{(n)}$ è omotopa all'identità tramite un'omotopia lipschitziana
$H : [0,1]\times \mathbb{R}^n \to \mathbb{R}^n$.

Dalla formula di omotopia per correnti si ottiene:
\[
T^{(n-1)} - T = \partial S^{(n)} + R^{(n)},
\]
dove:
\begin{itemize}
\item $S^{(n)}$ è una $(k+1)$-corrente associata all'omotopia,
\item $R^{(n)}$ è un termine di errore che dipende da $\partial T$.
\end{itemize}

Inoltre si ha la stima:
\[
\mathbf{M}(S^{(n)}) \le C\,2^{-N}\,\mathbf{M}(T),
\]
poiché lo spostamento indotto da $\pi^{(n)}$ è dell'ordine di $2^{-N}$.

\medskip

\noindent
\textbf{Passo 3: iterazione.}

Si ripete lo stesso procedimento sugli scheletri successivi.

Costruiamo mappe:
\[
\pi^{(n-1)}, \pi^{(n-2)}, \dots, \pi^{(k+1)}
\]
che deformano rispettivamente:
\[
(n-1)\text{-scheletro} \to (n-2)\text{-scheletro} \to \cdots \to k\text{-scheletro}.
\]

Definendo ricorsivamente:
\[
T^{(j-1)} := \pi^{(j)}_{\#} T^{(j)},
\]
si ottiene alla fine una corrente:
\[
P_N := T^{(k)}
\]
supportata sullo scheletro $k$-dimensionale della griglia.

\medskip

\noindent
\textbf{Passo 4: decomposizione finale.}

Sommando le relazioni di omotopia ottenute a ogni passo,
si ottiene:
\[
T = P_N + \partial S_N + R_N,
\]
dove:
\[
S_N := \sum_j S^{(j)}, \qquad R_N := \sum_j R^{(j)}.
\]

\medskip

\noindent
\textbf{Passo 5: stime di massa.}

Le stime seguono dal fatto che:
\begin{itemize}
\item ogni deformazione sposta i punti di una quantità $\lesssim 2^{-N}$,
\item le mappe coinvolte hanno costante lipschitziana controllata.
\end{itemize}

Quindi:
\[
\mathbf{M}(S_N) \le C\,2^{-N}\,\mathbf{M}(T),
\]
\[
\mathbf{M}(R_N) \le C\,2^{-N}\,\mathbf{M}(\partial T).
\]

\medskip

\noindent
\textbf{Passo 6: caso senza bordo.}

Se $\partial T = 0$, allora tutti i termini $R^{(j)}$ sono nulli,
e quindi:
\[
T = P_N + \partial S_N.
\]

\medskip

\noindent
Questo conclude la dimostrazione.
\end{proof}

\end{document}